% =============================================================================
%  Notas de clase — Sesión 19: Proyectos de múltiples archivos
%  Programación Avanzada — Universidad Panamericana
%  Compilar: pdflatex sesion19_proyectos_multiarchivo.tex
% =============================================================================
\documentclass[12pt, oneside]{article}
\usepackage{progavanzada}
\usepackage{array}

\begin{document}

\portada{Sesión 19}{Proyectos de múltiples archivos}{2026}

\tableofcontents
\newpage

% =============================================================================
\section{1976: el bug que no era un bug}
% =============================================================================

En 1976, Stuart Feldman trabajaba en los Laboratorios Bell cuando un colega
pasó horas depurando un programa que juró haber corregido.  El problema no
estaba en el código: el ejecutable había sido compilado desde una versión
\textit{antigua} del fuente.  Alguien había editado el archivo, pero nadie
había vuelto a compilar.

Feldman decidió que las computadoras debían rastrear qué archivos habían
cambiado y recompilar solo lo necesario.  En semanas escribió
\textbf{Make}, la herramienta que desde entonces acompaña a todos los
proyectos de software en Unix.  Casi cincuenta años después, todavía la
usamos.

\begin{tecnico}[Historia: Make y el origen de los gestores de construcción]
  Stuart Feldman inventó Make en 1976 en Bell Labs.  Su diseño refleja la
  filosofía Unix: una sola herramienta que hace una cosa bien y se compone
  con otras.
  \par\smallskip
  Décadas después surgieron alternativas: \textbf{CMake} (1999) para
  proyectos multiplataforma, \textbf{Ninja} (2012) para velocidad extrema
  en proyectos grandes (Chromium, LLVM), y \textbf{Bazel} (2015, Google)
  para repositorios monolíticos de millones de líneas.
  \par\smallskip
  \textit{Referencia:} Feldman, S.\,I. (1979). ``Make — A Program for
  Maintaining Computer Programs.'' \textit{Software: Practice and
  Experience}, 9(4), 255–265.
\end{tecnico}

% =============================================================================
\section{El problema: cuando un archivo no alcanza}
% =============================================================================

Empieza un proyecto con 50 líneas.  Un mes después tiene 500; un semestre
después, 5\,000.  Editar y recompilar un solo archivo gigante tiene varios
costos:

\begin{itemize}
  \item \textbf{Lectura:} buscar una función entre miles de líneas es lento.
  \item \textbf{Colaboración:} dos personas no pueden editar el mismo archivo
        sin conflictos de versión.
  \item \textbf{Compilación:} cambiar una función obliga a recompilar todo el
        programa, aunque el cambio no afecte al resto.
  \item \textbf{Reutilización:} no puedes usar tus funciones en otro proyecto
        sin copiar el archivo completo.
\end{itemize}

La solución es dividir el código en \textbf{módulos}: unidades independientes
que se pueden escribir, compilar y reutilizar por separado.

% =============================================================================
\section{Archivos de cabecera: el contrato entre módulos}
% =============================================================================

\subsection{Declaración versus definición}

Cada módulo existe en dos archivos:

\begin{description}
  \item[\texttt{modulo.h}] Contiene las \textbf{declaraciones}: le dice al
    compilador qué funciones y tipos \textit{existen} y cuál es su forma.
    No hay código ejecutable.
  \item[\texttt{modulo.cpp}] Contiene las \textbf{definiciones}: cómo
    funcionan esas funciones.  Incluye su propio \texttt{.h} para que el
    compilador verifique que la definición coincide con la declaración.
\end{description}

\textbf{funciones.h} — solo declaraciones:
\begin{verbatim}
void saludar(char*);
int  sumar(int, int);
int  restar(int, int);
\end{verbatim}

\textbf{funciones.cpp} — definiciones (incluye su propio \texttt{.h}):
\begin{verbatim}
#include "funciones.h"
void saludar(char* n) { cout << "Hola " << n << endl; }
int  sumar(int a, int b) { return a + b; }
int  restar(int a, int b) { return a - b; }
\end{verbatim}

La regla de C++ que gobierna esto se llama \textbf{ODR} (\textit{One
Definition Rule}): puedes \textit{declarar} algo cuantas veces quieras,
pero solo puedes \textit{definirlo} una vez en todo el programa.

\subsection{Guards de inclusión}

¿Qué pasa si \texttt{main.cpp} incluye \texttt{funciones.h} directamente
\textbf{y} también incluye \texttt{auxiliar.h}, que a su vez incluye
\texttt{funciones.h}?  El preprocesador copiaría el contenido de
\texttt{funciones.h} dos veces, produciendo un error de redefinición.

Los \textbf{guards de inclusión} lo evitan:

\begin{verbatim}
#ifndef FUNCIONES_H   // Si la macro no existe...
#define FUNCIONES_H   // ...crearla (primera inclusión)

void saludar(char*);
int  sumar(int, int);
int  restar(int, int);

#endif                // Fin del bloque protegido
\end{verbatim}

La segunda vez que el preprocesador encuentra \verb|#include "funciones.h"|,
la macro \texttt{FUNCIONES\_H} ya existe y salta todo el contenido.

\begin{consejo}
  La alternativa moderna es escribir \cpp{\#pragma once} como primera
  línea del \texttt{.h}.  Todos los compiladores actuales la soportan
  y es más simple que el guard manual.  Ambas formas son correctas.
\end{consejo}

\subsection{\texttt{<>} vs. \texttt{""}}

\begin{center}
\begin{tabular}{lll}
  \toprule
  \textbf{Sintaxis} & \textbf{Busca en} & \textbf{Usar para} \\
  \midrule
  \cpp{\#include <vector>}       & Rutas del sistema          & Librería estándar y paquetes instalados \\
  \cpp{\#include "funciones.h"}  & Directorio actual primero  & Tus propios archivos \\
  \bottomrule
\end{tabular}
\end{center}

Regla simple: si el archivo lo escribiste tú, usa \texttt{""}; si viene
con el compilador o una librería externa, usa \texttt{<>}.

% =============================================================================
\section{Compilación por etapas}
% =============================================================================

\subsection{Los dos pasos del compilador}

Cuando escribes \texttt{g++ main.cpp funciones.cpp -o programa}, en
realidad ocurren dos pasos distintos:

\begin{center}
\begin{tabular}{lll}
  \toprule
  \textbf{Paso} & \textbf{Entrada} & \textbf{Salida} \\
  \midrule
  1.\ Compilar  & \texttt{main.cpp}      & \texttt{main.o}      \\
                & \texttt{funciones.cpp} & \texttt{funciones.o} \\
  \midrule
  2.\ Enlazar   & \texttt{main.o} + \texttt{funciones.o} & \texttt{programa} \\
  \bottomrule
\end{tabular}
\end{center}

El \textbf{compilador} traduce cada \texttt{.cpp} a código máquina de
forma independiente, produciendo un \textbf{archivo objeto} (\texttt{.o}).
El \textbf{enlazador} (\textit{linker}) conecta esos objetos: cuando
\texttt{main.o} llama a \texttt{saludar}, el linker busca en
\texttt{funciones.o} dónde está esa función y resuelve la referencia.

\subsection{Comandos individuales}

\begin{verbatim}
# Paso 1: compilar cada .cpp por separado (genera .o, no enlaza)
g++ -c funciones.cpp -o funciones.o
g++ -c main.cpp       -o main.o

# Paso 2: enlazar todos los .o en un ejecutable
g++ main.o funciones.o -o programa
\end{verbatim}

\begin{recuerda}
  La bandera \texttt{-c} le dice a \texttt{g++} que solo compile
  (no enlace).  Sin \texttt{-c}, intenta compilar y enlazar todo lo que
  reciba.
\end{recuerda}

\subsection{Una línea de compilación que no para de crecer}

Un proyecto de dos archivos es manejable.  Pero imagina el proyecto de
personas que veremos al final de esta sesión: seis módulos, cada uno con
su \texttt{.h} y su \texttt{.cpp}.  La línea de compilación queda así:

\begin{verbatim}
g++ -std=c++17 -Wall -Wextra -o programa \
    main.cpp auxiliar.cpp obtener_datos.cpp \
    mate.cpp buscar.cpp
\end{verbatim}

Y eso es solo al inicio.  Cuando el proyecto crezca y necesites flags
adicionales, rutas de include o librerías externas:

\begin{verbatim}
g++ -std=c++17 -Wall -Wextra -O2 -I./include \
    -I./lib/json/include \
    main.cpp auxiliar.cpp obtener_datos.cpp    \
    mate.cpp buscar.cpp graficas.cpp reportes.cpp \
    -L./lib -ljson -lm -o programa
\end{verbatim}

Cada vez que cambies un archivo tienes que recordar y volver a escribir
todo eso.  Cada vez que agregues un \texttt{.cpp} tienes que buscarlo
en la línea y añadirlo en el lugar correcto.

Hay una solución mejor.

% =============================================================================
\section{Make: automatizar la compilación}
% =============================================================================

\subsection{Un primer Makefile: solo la línea de compilación}

Un \texttt{Makefile} es un archivo de texto que le dice a Make qué
comandos ejecutar para construir tu programa.  La forma más básica es
simplemente envolver la línea de compilación en una \textbf{regla}:

\begin{verbatim}
programa:
	g++ -std=c++17 -Wall -Wextra -o programa \
	    main.cpp auxiliar.cpp obtener_datos.cpp \
	    mate.cpp buscar.cpp
\end{verbatim}

Con esto, en lugar de recordar y teclear la línea larga, solo escribes:

\begin{verbatim}
make
\end{verbatim}

Make busca el archivo llamado \texttt{Makefile} en el directorio actual,
lee la primera regla y ejecuta el comando.

\begin{advertencia}
  El comando dentro de una regla debe comenzar con un carácter de
  \textbf{tabulación} (Tab), no con espacios.  Es el error más común al
  escribir Makefiles: la mayoría de editores muestran ambos igual, pero
  Make los distingue y falla si encuentra espacios.
\end{advertencia}

La estructura de cualquier regla es siempre la misma:

\begin{verbatim}
objetivo: dependencias
<TAB>   comando para construir el objetivo
\end{verbatim}

\textbf{objetivo} es el nombre del archivo que se quiere generar.
\textbf{dependencias} son los archivos que necesita para generarlo.
Por ahora las dejamos vacías; las completaremos pronto.

\subsection{Variables: el compilador y sus flags}

Repetir \texttt{g++} y las flags en cada regla es un problema: si cambias
de compilador o decides agregar un flag nuevo, tienes que modificar cada
línea.  Las \textbf{variables} de Make resuelven esto.

Primero extraemos el compilador:

\begin{verbatim}
CXX = g++

programa:
	$(CXX) -std=c++17 -Wall -Wextra -o programa \
	    main.cpp auxiliar.cpp obtener_datos.cpp \
	    mate.cpp buscar.cpp
\end{verbatim}

Luego los flags:

\begin{verbatim}
CXX      = g++
CXXFLAGS = -std=c++17 -Wall -Wextra

programa:
	$(CXX) $(CXXFLAGS) -o programa \
	    main.cpp auxiliar.cpp obtener_datos.cpp \
	    mate.cpp buscar.cpp
\end{verbatim}

Para usar el valor de una variable se escribe \texttt{\$(NOMBRE)}.
Make sustituye el valor antes de ejecutar el comando.

\subsection{Variables: la lista de archivos y el nombre del ejecutable}

El mismo principio aplica a la lista de fuentes y al nombre del ejecutable:

\begin{verbatim}
CXX      = g++
CXXFLAGS = -std=c++17 -Wall -Wextra
TARGET   = programa
SRCS     = main.cpp auxiliar.cpp obtener_datos.cpp \
           mate.cpp buscar.cpp

$(TARGET): $(SRCS)
	$(CXX) $(CXXFLAGS) $(SRCS) -o $(TARGET)
\end{verbatim}

Ahora agregar un módulo nuevo es editar solo la variable \texttt{SRCS}.
Cambiar el nombre del ejecutable es editar solo \texttt{TARGET}.

\begin{ejemplo}[Agregar un módulo nuevo]
  El proyecto crece y necesitas un módulo \texttt{graficas.cpp}.
  Con el Makefile anterior solo cambias una línea:
  \begin{verbatim}
SRCS = main.cpp auxiliar.cpp obtener_datos.cpp \
       mate.cpp buscar.cpp graficas.cpp
  \end{verbatim}
  La próxima vez que corras \texttt{make}, el comando completo ya
  incluirá \texttt{graficas.cpp} automáticamente.
\end{ejemplo}

\subsection{Dependencias y archivos objeto}

El Makefile anterior todavía tiene un problema: recompila \textit{todos}
los fuentes cada vez que cambias cualquier cosa.  Si modificas solo
\texttt{buscar.cpp}, Make vuelve a compilar \texttt{main.cpp},
\texttt{auxiliar.cpp} y los demás aunque no los tocaste.

La solución es declarar cada \texttt{.o} como objetivo separado, con sus
propias dependencias:

\begin{verbatim}
CXX      = g++
CXXFLAGS = -std=c++17 -Wall -Wextra
TARGET   = programa
OBJS     = main.o auxiliar.o obtener_datos.o mate.o buscar.o

$(TARGET): $(OBJS)
	$(CXX) $(OBJS) -o $(TARGET)

main.o: main.cpp auxiliar.h obtener_datos.h mate.h buscar.h
	$(CXX) $(CXXFLAGS) -c main.cpp

auxiliar.o: auxiliar.cpp auxiliar.h persona.h
	$(CXX) $(CXXFLAGS) -c auxiliar.cpp

obtener_datos.o: obtener_datos.cpp obtener_datos.h persona.h auxiliar.h
	$(CXX) $(CXXFLAGS) -c obtener_datos.cpp

mate.o: mate.cpp mate.h persona.h
	$(CXX) $(CXXFLAGS) -c mate.cpp

buscar.o: buscar.cpp buscar.h persona.h
	$(CXX) $(CXXFLAGS) -c buscar.cpp

clean:
	rm -f $(OBJS) $(TARGET)
\end{verbatim}

Las dependencias de cada \texttt{.o} incluyen los \texttt{.h} que ese
módulo utiliza, no solo su propio \texttt{.cpp}.  Si cambias
\texttt{persona.h} (por ejemplo, agregas un campo), Make recompila todos
los módulos que dependen de él.

\subsection{Cómo decide Make qué recompilar}

Make compara la \textbf{fecha de modificación} de cada objetivo con la de
sus dependencias.  Si alguna dependencia es más reciente que el objetivo,
lo reconstruye; si no, no hace nada.

\begin{verbatim}
# Solo cambias buscar.cpp:
$ make
g++ -std=c++17 -Wall -Wextra -c buscar.cpp      # solo este
g++ main.o auxiliar.o obtener_datos.o mate.o \
    buscar.o -o programa                         # reenlaza
\end{verbatim}

\texttt{main.o}, \texttt{auxiliar.o} y los demás no se tocan: sus fechas
son más recientes que sus fuentes.  En proyectos con cientos de archivos,
esto reduce el tiempo de compilación de minutos a segundos.

\subsection{Regla patrón: evitar la repetición}

Escribir una regla distinta para cada \texttt{.cpp} es tedioso.  Una
\textbf{regla patrón} lo generaliza:

\begin{verbatim}
# Cualquier .o se construye desde el .cpp con el mismo nombre
%.o: %.cpp
	$(CXX) $(CXXFLAGS) -c $< -o $@
\end{verbatim}

Las variables automáticas \texttt{\$<} y \texttt{\$@} son especiales de Make:

\begin{center}
\begin{tabular}{ll}
  \toprule
  \textbf{Variable} & \textbf{Significado} \\
  \midrule
  \texttt{\$@} & El objetivo actual (\texttt{buscar.o}) \\
  \texttt{\$<} & La primera dependencia (\texttt{buscar.cpp}) \\
  \texttt{\$\^{}} & Todas las dependencias \\
  \bottomrule
\end{tabular}
\end{center}

Con la regla patrón, el Makefile completo queda mucho más corto:

\begin{verbatim}
CXX      = g++
CXXFLAGS = -std=c++17 -Wall -Wextra
TARGET   = programa
SRCS     = main.cpp auxiliar.cpp obtener_datos.cpp mate.cpp buscar.cpp
OBJS     = $(SRCS:.cpp=.o)

$(TARGET): $(OBJS)
	$(CXX) $(OBJS) -o $(TARGET)

%.o: %.cpp
	$(CXX) $(CXXFLAGS) -c $< -o $@

clean:
	rm -f $(OBJS) $(TARGET)
\end{verbatim}

La notación \texttt{\$(SRCS:.cpp=.o)} le pide a Make que sustituya
\texttt{.cpp} por \texttt{.o} en cada nombre de \texttt{SRCS}: es
equivalente a escribir \texttt{OBJS = main.o auxiliar.o ...} a mano.

% =============================================================================
\section{CMake: multiplataforma y sin repetición}
% =============================================================================

\subsection{¿Por qué CMake si ya tenemos Make?}

Make funciona perfectamente en Linux y Mac.  Pero en Windows la sintaxis
de las terminales es diferente (\texttt{rm} no existe, las rutas usan
\texttt{\textbackslash} en lugar de \texttt{/}), y los Makefiles escritos
en Linux no funcionan ahí sin modificaciones.

\textbf{CMake} resuelve esto: describes tu proyecto una sola vez en un
archivo llamado \texttt{CMakeLists.txt} y CMake genera, en cada
plataforma, el sistema de compilación nativo.

\begin{center}
\begin{tabular}{ll}
  \toprule
  \textbf{Plataforma} & \textbf{CMake genera} \\
  \midrule
  Linux / Mac        & \texttt{Makefile} \\
  Windows            & Proyecto de Visual Studio (\texttt{.sln}) \\
  Cualquiera (opcional) & Archivos Ninja (más rápido que Make) \\
  \bottomrule
\end{tabular}
\end{center}

\subsection{Un \texttt{CMakeLists.txt} mínimo}

\begin{verbatim}
cmake_minimum_required(VERSION 3.10)

project(MiProyecto VERSION 1.0)

set(CMAKE_CXX_STANDARD 17)
set(CMAKE_CXX_STANDARD_REQUIRED True)

add_executable(programa
    main.cpp
    auxiliar.cpp
    obtener_datos.cpp
    mate.cpp
    buscar.cpp
)
\end{verbatim}

Es más verboso que el Makefile, pero no tiene ninguna sintaxis
dependiente del sistema operativo.  El mismo \texttt{CMakeLists.txt}
funciona en Linux, Mac y Windows sin cambios.

\subsection{Flujo de trabajo}

CMake separa el código fuente de los archivos de compilación mediante
un \textbf{build fuera de la fuente} (\textit{out-of-source build}):

\begin{verbatim}
mkdir build   # crear directorio de compilación
cd build
cmake ..      # generar el Makefile (lee CMakeLists.txt del directorio padre)
make          # compilar
./programa    # ejecutar
\end{verbatim}

Los archivos generados (\texttt{.o}, \texttt{Makefile}, etc.) quedan en
\texttt{build/} y \textbf{no se suben al repositorio}.  Para empezar de
cero basta con borrar la carpeta \texttt{build/} y repetir el proceso.

\begin{consejo}
  Agrega \texttt{build/} a tu \texttt{.gitignore} para que git nunca
  lo rastree.  Solo \texttt{CMakeLists.txt} y los archivos fuente deben
  vivir en el repositorio.
\end{consejo}

% =============================================================================
\section{Proyecto ejemplo: sistema de personas}
% =============================================================================

El directorio \texttt{materialInicial/sesion19/proyecto1/} tiene un
proyecto real con siete archivos.  La siguiente tabla muestra qué
responsabilidad tiene cada módulo:

\begin{center}
\begin{tabular}{lp{9cm}}
  \toprule
  \textbf{Módulo} & \textbf{Responsabilidad} \\
  \midrule
  \texttt{persona.h}             & Define la estructura \texttt{Persona} (nombre, apellidos, edad, salario).  Sin lógica. \\
  \texttt{auxiliar.h/.cpp}       & Genera nombres y apellidos aleatorios de listas predefinidas. \\
  \texttt{obtener\_datos.h/.cpp} & Construye un arreglo de \texttt{Persona} con datos generados. \\
  \texttt{mate.h/.cpp}           & Calcula promedios de edad y salario. \\
  \texttt{buscar.h/.cpp}         & Ordena y filtra arreglos de \texttt{Persona}. \\
  \texttt{main.cpp}              & Orquesta todo: genera datos, calcula y muestra resultados. \\
  \bottomrule
\end{tabular}
\end{center}

El grafo de inclusión muestra por qué \texttt{persona.h} está en el
centro: todos los módulos trabajan con el tipo \texttt{Persona}, así que
todos lo incluyen.

\begin{verbatim}
              main.cpp
             /   |   \
     buscar.h  mate.h  obtener_datos.h
        |         |          |
    persona.h persona.h  persona.h
                            |
                        auxiliar.h
\end{verbatim}

Los guards de inclusión en \texttt{persona.h} garantizan que, aunque
aparezca cuatro veces en la cadena de inclusiones, el preprocesador solo
la procese una vez.

% =============================================================================
\section{Ejercicios}
% =============================================================================

\begin{enumerate}

  \item Dado el proyecto mínimo (\texttt{main.cpp}, \texttt{funciones.h},
        \texttt{funciones.cpp}), escribe el \texttt{Makefile} más simple
        posible: una sola regla con la línea de compilación directa (sin
        variables).  Luego evoluciona ese Makefile agregando
        \texttt{CXX}, \texttt{CXXFLAGS}, \texttt{TARGET} y \texttt{SRCS}
        paso a paso.

  \item Con el Makefile final de la sección anterior (usando
        \texttt{OBJS} y la regla patrón), modifica \texttt{funciones.cpp}
        y ejecuta \texttt{make}.  ¿Qué archivos recompila y cuáles omite?
        Repite el experimento modificando solo \texttt{funciones.h}.
        ¿Cambia el resultado?

  \item ¿Qué ocurre si defines una función dos veces en el mismo
        \texttt{.cpp}?  ¿Y si la defines en dos \texttt{.cpp} distintos
        que se enlazan juntos?  Prueba ambos casos y explica el mensaje
        de error en cada uno.

  \item Crea un nuevo módulo \texttt{estadisticas.h} /
        \texttt{estadisticas.cpp} para el proyecto de personas con la
        función \cpp{double salario\_maximo(Persona arr[], int n)}.
        Incluye el guard de inclusión correcto, agrégalo a \texttt{SRCS}
        en el Makefile, y escribe la línea equivalente en
        \texttt{CMakeLists.txt}.

  \item ¿Cuál es la diferencia entre \texttt{make} y \texttt{make clean}?
        ¿Cuándo conviene ejecutar \texttt{clean} antes de recompilar?

\end{enumerate}

% =============================================================================
\section*{Resumen}
% =============================================================================

\begin{itemize}
  \item Dividir el código en módulos (\texttt{.h} + \texttt{.cpp}) mejora
        la legibilidad, la colaboración y la velocidad de compilación.
  \item El \textbf{archivo de cabecera} contiene declaraciones (el
        \textit{qué}); el \textbf{archivo de implementación} contiene
        definiciones (el \textit{cómo}).
  \item Los \textbf{guards de inclusión} (\cpp{\#ifndef}/\cpp{\#define}/\cpp{\#endif}
        o \cpp{\#pragma once}) evitan que el mismo \texttt{.h} se procese
        más de una vez.
  \item La compilación tiene dos pasos: \textbf{compilar}
        (\texttt{.cpp} → \texttt{.o}) y \textbf{enlazar}
        (\texttt{.o} → ejecutable).
  \item Separar los pasos permite recompilar solo los módulos que cambiaron.
  \item Un \textbf{Makefile} puede empezar tan simple como una sola regla
        con la línea de compilación, y evolucionar con variables
        (\texttt{CXX}, \texttt{CXXFLAGS}, \texttt{SRCS}, \texttt{OBJS})
        conforme el proyecto crece.
  \item \textbf{Make} automatiza la recompilación: compara fechas de
        modificación y reconstruye solo lo que cambió.
  \item \textbf{CMake} genera el sistema de compilación adecuado para cada
        plataforma a partir de un \texttt{CMakeLists.txt} portátil.
\end{itemize}

\end{document}
